% Chapter 11: CI/CD, Containerization & Kubernetes Deployment

\begin{frame}[plain]
    \vfill
    \centering
    {\footnotesize \textbf{\textcolor{deepblue}{\textsc{TaxonVision-MLOps Architecture}}}} \\[0.35cm]
    {\large \textbf{\textcolor{signalorange}{CHAPTER 11}}} \\[0.25cm]
    {\huge \textbf{CI/CD \& Kubernetes Deployment}} \\[0.35cm]
    {\normalsize \textcolor{black!70}{Declarative Jenkins Pipelines, GHCR Containers \& Bare-Metal k3s Orchestration}} \\[0.5cm]
    \textcolor{deepblue}{\rule{0.55\linewidth}{1.2pt}} \\[0.5cm]
    \statusdone
    \vfill
    \script{In Chapter 11, we walk through our continuous delivery pipeline: bridging GitHub Actions with on-premise Jenkins, building versioned OCI containers, and orchestrating zero-downtime rollouts in k3s.}
    \qa{Why use a dual CI/CD system with GitHub Actions and Jenkins?}{GitHub Actions handles lightweight PR verification in the cloud, while private on-premise Jenkins builds heavy GPU images and deploys to our local bare-metal k3s cluster safely.}
\end{frame}

\begin{frame}{CI/CD Engineering}{End-to-End Workflow: The Dual-Pipeline CI/CD Bridge}
    \begin{block}{Security \& Network Boundary Architecture}
        \begin{itemize}\tiny
            \item \textbf{Public Cloud Enclave (GitHub Actions):} Pull request validation (Ruff linting, Mypy strict typing, complexipy cognitive complexity $\le 15$, open license audit, OKF validation). Zero access to bare-metal hardware or production secrets.
            \item \textbf{The Bridge:} PR merge to \texttt{develop} or \texttt{main} $\to$ Signed commit $\to$ Cloudflare Zero Trust outbound-authenticated webhook to private Jenkins controller on \texttt{hive-mind}.
            \item \textbf{Private On-Premise Enclave (Jenkins \& k3s):} Secret injection via Linux Secret Service / unified HMAC tokens. Dedicated GPU training and k3s deployment with zero inbound firewall ports opened.
        \end{itemize}
    \end{block}

    \vspace{0.1cm}
    \statusdone

    \script{Here is an enterprise DevOps pattern every student must learn: how do you connect cloud GitHub to private bare-metal servers securely? You never allow GitHub Actions to SSH directly into your private GPU cluster. Instead, GitHub Actions runs fast cloud checks on public runners. Once code is approved, a Cloudflare Zero Trust webhook securely bridges the gap. Jenkins running inside our private network pulls the signed commit, injects secrets, and executes the GPU deployment locally with zero inbound ports exposed.}
    \qa{Why not use GitHub Self-Hosted Runners directly on hive-mind?}{Self-hosted runners can execute untrusted code from malicious PRs directly on your machine. Our pipeline ensures only trusted code reaches bare-metal.}
    \qa{How does the Cloudflare Zero Trust webhook enhance security?}{It acts as an outbound-authenticated tunnel, meaning our private Jenkins controller does not need a public IP or open inbound firewall ports to receive the merge event.}
\end{frame}

\begin{frame}{Cloud Quality Gates}{GitHub Actions Workflow Suite: Isolation \& Public PR Gating}
    \begin{block}{Core Workflow: \texttt{.github/workflows/ci.yml} (PR \& Push Matrix)}
        \begin{itemize}\tiny
            \item \textbf{Job 1 (Code, Security \& Type Quality):} Ubuntu-latest CPU \texttt{ci-dev} environment; DagsHub DVC model checkpoint hydration; pre-commit hooks (\texttt{SKIP=enforce-author-identity,openwiki-sync}); OKF v0.2 validation; two-pass pytest (Pass 1: Unit + Coverage $\ge 80\%$, Pass 2: Integration + Conformal Invariants).
            \item \textbf{Job 2 (Complexipy Cognitive Complexity):} SARIF scan enforcing cognitive complexity $\le 15$; results uploaded via CodeQL \texttt{upload-sarif@v4} directly to GitHub Code Scanning alerts.
            \item \textbf{Job 3 (Architecture \& Line Count Profile):} Boyter SCC container scan (\texttt{ghcr.io/boyter/scc:master}) analyzing LOC and codebase composition.
            \item \textbf{Job 4 (Documentation Strict Build):} Generates OKF D3 knowledge graph, compiles Zensical MkDocs with \texttt{-{}-strict}, publishes to GitHub Pages on \texttt{main} merge.
            \item \textbf{Job 5 (Build \& Publish Serving Container):} Docker Buildx builds production image and publishes to GHCR tagged with commit SHA and \texttt{latest}.
        \end{itemize}
    \end{block}

    \begin{block}{Auxiliary Workflows: Automated Repository Maintenance}
        \begin{itemize}\tiny
            \item \textbf{Scheduled Evidence Indexing (\texttt{openwiki-update.yml}):} Scheduled weekly cron (Monday 08:00 UTC) running Gemini 3.8 Flash to index claims, architecture, and tests into \texttt{openwiki/}.
            \item \textbf{Semantic Tag Release Pipeline (\texttt{release.yaml}):} Triggers on \texttt{v*} tags; verifies version alignment between \texttt{pyproject.toml} and \texttt{taxon\_vision.\_\_version\_\_}, packages wheels, and drafts GitHub Releases.
        \end{itemize}
    \end{block}

    \vspace{0.1cm}
    \statusdone

    \script{On our public cloud perimeter, GitHub Actions serves as the first line of defense. The primary CI workflow executes five decoupled jobs: static code quality with strict eighty percent coverage enforcement, SARIF-based cognitive complexity scanning with Complexipy, code profiling, Zensical documentation builds deployed to GitHub Pages, and container image publishing to GHCR. Two auxiliary workflows handle automated OpenWiki knowledge graphs every Monday and semantic release packaging whenever a tag is pushed. Crucially, public cloud runners never touch private bare-metal servers or production credentials.}
    \qa{Why split the test suite into two passes in ci.yml?}{Pass 1 runs fast unit tests and fails immediately if total branch coverage falls below eighty percent. Pass 2 executes heavier integration and conformal invariant tests in fail-fast mode.}
    \qa{How does the complexipy job integrate with GitHub security features?}{It generates an industry-standard SARIF report, uploaded via the CodeQL action directly into the GitHub Security tab as automated code-quality advisories.}
\end{frame}

\begin{frame}{Jenkins Automation}{Jenkinsfile Execution: Step-by-Step Pipeline Trace}
    \begin{block}{Six Declarative Pipeline Stages}
        \begin{itemize}\scriptsize
            \item \textbf{Stage 1 (Checkout \& Secrets):} Ephemeral workspace hydration; fetches credentials from Linux Secret Service / DagsHub token.
            \item \textbf{Stage 2 (Environment Hydration):} \texttt{pixi install {-}{-}frozen} with NVMe Rattler cache mount (\texttt{/root/.cache/rattler/cache}) avoiding redundant package downloads.
            \item \textbf{Stage 3 (DVC Pull):} Fast download of Parquet manifests and embeddings via \texttt{s3://dvc} Content-Addressed Storage.
            \item \textbf{Stage 4 (GPU Training \& Optuna):} Strict 5.0 GB VRAM boundary, 150-minute stage timeout (pipeline 240m), \texttt{PYTHONUNBUFFERED=1} unbuffered live logging to console.
            \item \textbf{Stage 5 (ONNX Export \& INT8 PTQ):} Static spatial $3 \times 224 \times 224$, dynamic batch \texttt{\{0: "batch\_size"\}}, QDQ operator insertion.
            \item \textbf{Stage 6 (Artifact Push \& Promotion):} \texttt{dvc push} of final weights and GHCR OCI container tag promotion.
        \end{itemize}
    \end{block}

    \vspace{0.2cm}
    \statusdone

    \script{Once the webhook triggers Jenkins, the declarative Jenkinsfile executes our six-stage pipeline. It starts by hydrating the workspace and injecting credentials. Stage 2 is crucial: by mounting the NVMe Rattler cache, Pixi avoids re-downloading hundreds of megabytes of Conda packages. The pipeline then pulls DVC data, runs Optuna tuning with strict VRAM limits, exports the ONNX models with dynamic batching, and finally pushes the built artifacts to GHCR. Unbuffered logging ensures we see real-time output even during long 150-minute GPU training stages.}
    \qa{Why do we need a 5.0 GB VRAM boundary in Stage 4?}{Our bare-metal RTX 5070 Ti uses NVIDIA time-slicing to share its 16 GB VRAM across different workloads. Enforcing a strict 5.0 GB ceiling for the Jenkins training pod prevents CI pipelines from crashing the live production FastAPI inference server.}
    \qa{Why is \texttt{PYTHONUNBUFFERED=1} important for CI logs?}{By default, Python buffers stdout, which means logs might not appear in the Jenkins console for minutes during heavy ML training loops. Unbuffered logging flushes output instantly, allowing developers to monitor progress in real time.}
\end{frame}

\begin{frame}{On-Premise Automation}{Declarative Jenkinsfile: Parameters, Agent Pod \& 7-Stage DAG}
    \vspace{-0.15cm}
    \begin{block}{Declarative Pipeline Envelope \& Agent Specification}
        \begin{itemize}\tiny
            \item \textbf{Pipeline Parameters:} \texttt{RETRAIN\_MODEL} (boolean flag to force GPU retraining), \texttt{TUNE\_HYPERPARAMETERS} (boolean to trigger Optuna search), \texttt{EXTRACTOR} (string backbone selector, default \texttt{dinov3}).
            \item \textbf{Kubernetes Ephemeral Agent:} Dynamically provisioned in \texttt{jenkins-ci} namespace; runs \texttt{ghcr.io/prefix-dev/pixi:latest}; mounts \texttt{emptyDir} \texttt{rattler-cache} at \texttt{/root/.cache/rattler/cache} for fast package caching.
            \item \textbf{Build Governance \& Credentials:} \texttt{buildDiscarder} retains 15 builds/14 days; global 240m pipeline timeout; \texttt{DAGSHUB\_TOKEN} credential binding; \texttt{PYTHONUNBUFFERED=1} and thread limits (\texttt{OMP}, \texttt{MKL}, \texttt{OPENBLAS=4}).
        \end{itemize}
    \end{block}

    \begin{block}{7-Stage Enterprise Execution Architecture}
        \begin{itemize}\tiny
            \item \textbf{Stage 1 (Prepare Toolchain):} Rattler cache mount, git/kubectl check, and DVC pull checkpoint hydration.
            \item \textbf{Stage 2 (Static Quality \& Invariants):} Parallel execution of Ruff lint/format, Mypy strict typing, and license/OKF validation.
            \item \textbf{Stage 3--4 (Unit \& Integration Tests):} Pytest unit suite; integration and conformal coverage invariants with coverage append ($\ge 80\%$).
            \item \textbf{Stage 5 (Documentation Strict Build):} Automated OKF graph visualization and Zensical strict documentation build.
            \item \textbf{Stage 6 (Model Training GPU):} Conditional on \texttt{develop}/\texttt{main} with changeset in \texttt{data/**}, \texttt{models/**}, \texttt{pipelines/**} or \texttt{RETRAIN\_MODEL}; swaps to GPU agent pod with 150m stage timeout.
            \item \textbf{Stage 7 (Deploy to Kubernetes):} Rolling deployment to \texttt{taxon-vision} namespace via \texttt{kubectl apply -f deploy/k8s/api-deployment.yaml} and rollout status verification.
        \end{itemize}
    \end{block}

    \vspace{0.02cm}
    \statusdone

    \script{When code merges to develop or main, on-premise Jenkins takes over via declarative pipeline with dynamic parameters. It runs on an ephemeral k3s agent, mounting an emptyDir Rattler cache for fast dependency hydration across seven stages: toolchain setup, parallel quality audits, unit tests, conformal invariants, documentation builds, GPU training in an isolated pod, and rolling production deployment.}
    \qa{Why mount an emptyDir volume for the Rattler package cache in the agent?}{Pixi uses the Rattler package solver and cache. Without an ephemeral cache mount, every build would re-download Conda packages from scratch.}
    \qa{How does Jenkins ensure training only runs when relevant code changes?}{Stage 6 uses declarative changeset conditions matching data, models, or pipelines paths, or the RETRAIN\_MODEL parameter.}
\end{frame}

\begin{frame}{Ephemeral GPU Runner}{K3s Pod Specification \& \texttt{jenkins\_gpu\_train.sh} Trace}
    \begin{block}{GPU Pod Specification (\texttt{deploy/k8s/jenkins-gpu-pod.yaml})}
        \begin{itemize}\tiny
            \item \textbf{Cluster \& Namespace Isolation:} Deployed to \texttt{jenkins-ci} namespace under dedicated \texttt{jenkins-agent} ServiceAccount.
            \item \textbf{Resource Quotas \& Time-Slicing:} Requests $2$ GiB RAM / $1000$m CPU; limits $16$ GiB RAM, $8$ CPU cores, and \texttt{nvidia.com/gpu: "1"} (1 virtual slice of RTX 5070 Ti).
            \item \textbf{Persistent Host Storage Mounts:} Rattler \texttt{emptyDir} cache; hostPath mounts for HuggingFace model weights (\texttt{/var/lib/taxon-vision/hf-cache}), Optuna SQLite database (\texttt{/var/lib/taxon-vision/optuna}), and embedding tensors (\texttt{/var/lib/taxon-vision/embeddings}).
        \end{itemize}
    \end{block}

    \begin{block}{Execution Lifecycle Trace (\texttt{scripts/ci/jenkins\_gpu\_train.sh})}
        \begin{itemize}\tiny
            \item \textbf{1. Bootstrap \& Credentials:} Installs git/kubectl via Pixi global; configures local SeaweedFS S3 credentials (\texttt{taxonadmin}).
            \item \textbf{2. Embedding Cache Priming:} Checks \texttt{/var/lib/taxon-vision/embeddings/*.pt} on host or pulls from cluster \texttt{local-seaweedfs} to bypass expensive feature extraction.
            \item \textbf{3. Pipeline Execution:} Invokes \texttt{dvc repro} with hyperparameter overrides (\texttt{-S training.tune=true}, \texttt{-S model.extractor=...}).
            \item \textbf{4. Multi-Tier Cache Synchronization:} Copies newly generated \texttt{.pt} tensors to host tier and pushes Epistemic Diamond to \texttt{local-seaweedfs}.
            \item \textbf{5. DagsHub Remote Push \& Lineage Lock:} Pushes \texttt{head.pt} and \texttt{model\_int8.onnx} to DagsHub S3 remote (\texttt{origin}); auto-commits \texttt{dvc.lock}, \texttt{conformal\_config.yaml}, and \texttt{reports/*.json} with \texttt{[skip ci]}.
        \end{itemize}
    \end{block}

    \vspace{0.1cm}
    \statusdone

    \script{For the GPU training stage, Jenkins dynamically schedules an ephemeral pod defined in jenkins-gpu-pod.yaml. This pod requests one virtual time-sliced GPU slice and mounts host directories for Hugging Face models, Optuna study databases, and pre-extracted embeddings. The script jenkins\_gpu\_train.sh bootstraps the environment, primes embeddings from the local host tier or SeaweedFS, runs dvc repro with any requested Optuna overrides, synchronizes fresh checkpoints across local and remote S3 tiers, and finally commits the updated dvc.lock and evaluation reports back to Git with a skip-ci annotation.}
    \qa{Why do we prime embeddings from a host directory before running dvc repro?}{Feature extraction over large image manifests is computationally expensive. If embeddings have not changed, copying pre-computed tensors from the host SSD saves hours of GPU compute.}
    \qa{Why does jenkins\_gpu\_train.sh push to both local SeaweedFS and DagsHub?}{Local SeaweedFS provides ultra-fast in-cluster caching for low-latency pod restarts, while DagsHub acts as our immutable external cloud backup and team collaboration hub.}
\end{frame}

\begin{frame}{Code Quality}{Cognitive Complexity Gates: complexipy $\le 15$ \& Strict Typing}
    \begin{block}{Automated Quality \& Complexity Enforcers}
        \begin{itemize}\small
            \item \textbf{The ``Research Spaghetti'' Problem:} ML codebases often decay into deeply nested loops, global states, and untyped tensors.
            \item \textbf{Strict Automated Gates:} TaxonVision enforces code quality through Ruff (formatting/linting) and Mypy (static typing).
            \item \textbf{Cognitive Complexity Limit:} \texttt{complexipy} restricts functions to a cognitive complexity of $\le 15$.
            \item \textbf{Refactoring Mandate:} Exceeding the complexity limit fails the pre-commit gate, forcing modular architecture and clean abstraction.
        \end{itemize}
    \end{block}

    \vspace{0.2cm}
    \statusdone

    \script{Why do so many data science projects fail when moving to production? Code complexity. In Jupyter notebooks, people write 300-line functions with 6 nested loops and zero type annotations. In TaxonVision, our CI gate uses \texttt{complexipy} to measure cognitive complexity. If any function scores higher than 15, the build fails instantly. This forces engineers to write modular, single-responsibility functions that are easy to unit test and maintain.}
    \qa{What was an example of a refactoring forced by this complexity rule?}{In \texttt{runner.py}, our embedding split preparation and index mapping functions originally scored above 20 due to combined dictionary parsing and tensor slicing. We decomposed them into dedicated helper methods, bringing complexity down to 6.}
\end{frame}

\begin{frame}{Container \& Cluster Topology}{Container Architecture: Multi-Stage Dockerfile vs. k3s}
    \begin{block}{Multi-Stage Container \& Deployment Spectrum}
        \begin{itemize}\footnotesize
            \item \textbf{Multi-Stage Dockerfile Optimization:}
            \begin{itemize}\scriptsize
                \item \textbf{Stage 1 (\texttt{builder}):} Ubuntu 24.04 with full Pixi toolchain, compilers, and Rattler build cache.
                \item \textbf{Stage 2 (\texttt{runtime}):} Stripped minimal Ubuntu image copying solely the compiled environment (\texttt{/app/.pixi/envs/default}) and application source. Final image size $< 1.5$ GB (reduced from $> 10$ GB).
            \end{itemize}
            \item \textbf{Local Development (\texttt{flagship}):} \texttt{deploy/compose/docker-compose.local.yaml} orchestration with local MinIO, mock DagsHub S3, and Prometheus scraper for zero-cloud offline debugging.
            \item \textbf{Bare-Metal Production (\texttt{hive-mind}):} Declarative Kubernetes (k3s) manifests with Traefik ingress controller, zero-downtime rolling updates, and NVIDIA GPU time-slicing daemonsets.
        \end{itemize}
    \end{block}

    \vspace{0.1cm}
    \statusdone

    \script{Let us examine how our containers transition from local development to production. We employ a multi-stage Dockerfile that compiles everything in a heavy builder stage, but only copies the minimal compiled environment into the runtime stage, shrinking our image from 10 gigabytes to under 1.5 gigabytes. Locally, developers use Docker Compose to mock cloud services like MinIO and DagsHub for offline debugging. In production on our bare-metal server, we orchestrate these same containers using Kubernetes manifests.}
    \qa{How does the multi-stage build securely handle credentials during the build process?}{By building in Stage 1 and copying compiled artifacts to Stage 2, any build-time secrets or SSH keys are absent from the final runtime image layer.}
    \qa{Why use Docker Compose locally but k3s in production?}{Docker Compose offers rapid iteration for local debugging. Production k3s provides robust self-healing, rolling deployments, Traefik ingress, and secure namespace isolation.}
\end{frame}

\begin{frame}{Deployment Strategy}{Zero-Downtime Rolling Updates \& Health Probes}
    \begin{block}{Zero-Downtime Rolling Update Topology}
        \begin{itemize}\tiny
            \item \textbf{The Downtime Challenge:} Updating an inference service must avoid dropping HTTP connections or returning 502 Bad Gateway errors.
            \item \textbf{Health Probes Configuration:} Kubernetes Deployment manifests define readiness and liveness probes targeting the FastAPI \texttt{/health} endpoint.
            \item \textbf{Model Warmup Sequence:} Kubernetes waits for model weights to load into GPU memory and warm up before declaring readiness.
            \item \textbf{Seamless Cutover:} Ingress traffic is diverted to new pods only when fully ready, and older replica pods are safely terminated.
        \end{itemize}
    \end{block}

    \vspace{0.2cm}
    \statusdone

    \script{Imagine updating your model in production: if you kill the old container before the new one finishes loading its 500-megabyte neural network into GPU memory, all incoming user requests fail with 500 errors. We configure Kubernetes rolling updates with readiness probes. The ingress keeps sending traffic to the old model until the new model reports that its weights are warmed up and ready. The switch happens seamlessly with zero dropped requests.}
    \qa{What is the difference between a liveness probe and a readiness probe?}{A liveness probe checks if the process is alive; if it fails, Kubernetes restarts the pod. A readiness probe checks if the pod is ready to accept HTTP traffic; if it fails, traffic is diverted without killing the container.}
\end{frame}

\begin{frame}{Cluster Topology}{Kubernetes Manifest Encyclopedia: RBAC, Storage \& Time-Slicing}
    \vspace{-0.15cm}
    \begin{block}{API Deployment \& Cluster Ingress (\texttt{api-deployment.yaml})}
        \begin{itemize}\tiny
            \item \textbf{Deployment Strategy:} 2 replicas with \texttt{RollingUpdate} (\texttt{maxSurge: 1}, \texttt{maxUnavailable: 0}) ensuring zero dropped connections.
            \item \textbf{Security Context:} Runs non-root (\texttt{runAsUser: 10001}, \texttt{runAsGroup: 10001}); read-only volume mount for checkpoints (\texttt{/app/models/checkpoints}).
            \item \textbf{Health Probing:} HTTP Liveness (\texttt{/health}, 15s delay, 20s period) and Readiness (\texttt{/health/ready}, 5s delay, 5s period, 2 failure threshold); exposed via NodePort service (\texttt{30080}).
        \end{itemize}
    \end{block}

    \begin{block}{RBAC, S3 Storage \& GPU Time-Slicing Subsystems}
        \begin{itemize}\tiny
            \item \textbf{Namespace Isolation \& RBAC (\texttt{jenkins-agent-rbac.yaml}):} \texttt{jenkins-ci} and \texttt{taxon-vision} namespaces; ServiceAccount \texttt{jenkins-agent} granted pod lifecycle rights in CI and deployment rollout rights in production.
            \item \textbf{In-Cluster S3 \& Optuna (\texttt{mlops-storage-and-optuna.yaml}):} SeaweedFS S3 server (\texttt{port 8333}) configured via \texttt{s3.json} (\texttt{taxon-operator} admin keys); persistent volumes for Optuna SQLite study DB and raw parquet datasets.
            \item \textbf{GPU Multiplexing (\texttt{nvidia-time-slicing-config.yaml}):} NVIDIA Device Plugin ConfigMap partitioning the single physical RTX 5070 Ti (16 GB) into 4 virtual slices (\texttt{replicas: 4}), allowing concurrent CI, serving, and training.
        \end{itemize}
    \end{block}

    \vspace{0.02cm}
    \statusdone

    \script{Our Kubernetes deployment on bare-metal k3s is governed by declarative manifests in deploy/k8s. The API runs two replicas with rolling update semantics, non-root user isolation, and liveness/readiness probes. RBAC policies enforce boundaries between jenkins-ci and taxon-vision. SeaweedFS provides in-cluster S3 storage for DVC and Optuna. NVIDIA time-slicing partitions our RTX 5070 Ti into four virtual GPU slices, preventing device deadlocks.}
    \qa{How does the readiness probe prevent 502 errors during rollouts?}{The readiness probe queries /health/ready, confirming ONNX Runtime sessions are fully initialized. Traefik routes zero traffic to new pods until it returns HTTP 200.}
    \qa{Why is NVIDIA time-slicing essential on a single-GPU bare-metal node?}{Consumer RTX GPUs lack hardware MIG. Without time-slicing, the first pod requesting a GPU locks the entire card exclusively, leaving other pods stuck in Pending.}
\end{frame}

\begin{frame}[fragile]{Pipeline Architecture Visualized}{Dual CI/CD Lifecycle Diagram}
    \centering
    \providecolor{TaxonPrimary}{RGB}{34, 139, 34}
    \providecolor{TaxonSecondary}{RGB}{0, 51, 102}
    \providecolor{TaxonDark}{RGB}{40, 40, 40}
    \providecolor{TaxonAccent}{RGB}{220, 120, 0}

    \resizebox{!}{0.58\textheight}{
    \begin{tikzpicture}[
        auto,
        block/.style={rectangle, draw=TaxonDark, thick, fill=white, text width=5.5em, text centered, rounded corners, minimum height=3em, font=\scriptsize},
        cloud/.style={rectangle, draw=TaxonSecondary, thick, fill=TaxonSecondary!10, text width=5.5em, text centered, rounded corners, minimum height=3em, font=\scriptsize},
        metal/.style={rectangle, draw=TaxonPrimary, thick, fill=TaxonPrimary!10, text width=5.5em, text centered, rounded corners, minimum height=3em, font=\scriptsize},
        line/.style={draw, thick, -latex, shorten >=2pt},
        label/.style={font=\tiny, text=TaxonDark, align=center}
    ]
        % Nodes
        \node [block] (pr) {PR Created};
        \node [cloud, right=1.2cm of pr] (gh) {GitHub Actions\\(Cloud Runners)};
        \node [block, right=1.2cm of gh] (merge) {PR Merge};
        \node [metal, below=1.2cm of merge] (webhook) {Webhook to\\ \texttt{hive-mind}};
        \node [metal, left=1.2cm of webhook] (jenkins) {k3s Jenkins\\Ephemeral Pod};
        \node [metal, left=1.2cm of jenkins] (build) {Build OCI\\Image};
        \node [cloud, below=1.2cm of build] (ghcr) {Push to\\GHCR};
        \node [metal, right=1.2cm of ghcr] (deploy) {k3s Rolling\\Update};

        % Connections
        \path [line] (pr) -- (gh);
        \path [line] (gh) -- node [label, above] {Ruff, Mypy,\\complexipy,\\audit} (merge);
        \path [line] (merge) -- (webhook);
        \path [line] (webhook) -- (jenkins);
        \path [line] (jenkins) -- (build);
        \path [line] (build) -- (ghcr);
        \path [line] (ghcr) -- node [label, above] {Pull Image} (deploy);
    \end{tikzpicture}
    }

    \script{This diagram visualizes our asymmetric pipeline. Notice the top half is the public cloud—GitHub Actions running lightweight static analysis. The bottom half is our secure bare-metal enclave. The only bridge between them is the webhook triggered on PR merge, ensuring that unvetted code never executes on our primary GPU hardware. Jenkins then builds the container, pushes it back to the cloud registry, and manages the local Kubernetes deployment.}
    \qa{Why does Jenkins build the image instead of GitHub Actions?}{Building a production ML container often requires compiling custom CUDA kernels or resolving large dependencies. Doing this on GitHub's free runners is slow and resource-constrained. Our bare-metal machine leverages local NVMe caches and powerful CPUs for rapid builds.}
\end{frame}

\begin{frame}[fragile]{Rolling Update Visualized}{Zero-Downtime Traefik Ingress Traffic Switching}
    \centering
    \providecolor{TaxonPrimary}{RGB}{34, 139, 34}
    \providecolor{TaxonSecondary}{RGB}{0, 51, 102}
    \providecolor{TaxonDark}{RGB}{40, 40, 40}
    \providecolor{TaxonAccent}{RGB}{220, 120, 0}
    \providecolor{PodGreen}{RGB}{200, 240, 200}
    \providecolor{PodRed}{RGB}{255, 200, 200}

    \resizebox{!}{0.62\textheight}{
    \begin{tikzpicture}[
        >=stealth,
        auto,
        node distance=2.5cm,
        ingress/.style={diamond, draw=TaxonSecondary, thick, fill=TaxonSecondary!20, text width=4em, text centered, minimum height=3em, font=\scriptsize\bfseries},
        podold/.style={rectangle, draw=TaxonDark, thick, fill=PodRed, text width=6em, text centered, rounded corners, minimum height=3em, font=\scriptsize},
        podnew/.style={rectangle, draw=TaxonPrimary, thick, fill=PodGreen, text width=6em, text centered, rounded corners, minimum height=3em, font=\scriptsize},
        textnode/.style={font=\scriptsize, align=center}
    ]
        % Ingress Node
        \node [ingress] (traefik) {Traefik\\Ingress};

        % V1 Pods (Old)
        \node [podold, right=3.5cm of traefik, yshift=1.5cm] (pod1) {Pod v1\\(Terminating)};
        \node [podold, right=3.5cm of traefik, yshift=0cm] (pod2) {Pod v1\\(Active)};

        % V2 Pods (New)
        \node [podnew, right=3.5cm of traefik, yshift=-1.5cm] (pod3) {Pod v2\\(Warming Up)};
        \node [podnew, right=3.5cm of traefik, yshift=-3cm] (pod4) {Pod v2\\(Ready)};

        % Traffic Routes
        \draw [->, thick, dashed, draw=gray] (traefik) -- node[above, font=\tiny, pos=0.55] {\textcolor{herbired}{\textbf{[Traffic Drained]}}} (pod1);
        \draw [->, thick, draw=TaxonSecondary] (traefik) -- node[above, font=\tiny, pos=0.55] {Live Traffic} (pod2);
        \draw [->, thick, dashed, draw=gray] (traefik) -- node[above, sloped, font=\tiny, pos=0.55] {Probe Pending} (pod3);
        \draw [->, thick, draw=TaxonPrimary] (traefik) -- node[above, sloped, font=\tiny, pos=0.55] {Live Traffic} (pod4);

        % Health Checks
        \node [textnode, right=1cm of pod3] (hc1) {Readiness Check:\\ \texttt{/health} $\rightarrow$ 503};
        \draw [->, dotted, thick] (pod3) -- (hc1);

        \node [textnode, right=1cm of pod4] (hc2) {Readiness Check:\\ \texttt{/health} $\rightarrow$ 200};
        \draw [->, dotted, thick] (pod4) -- (hc2);

    \end{tikzpicture}
    }

    \script{This diagram shows the transition state during a deployment. Traefik is the ingress controller mapping external requests. When we deploy Version 2, the new pod initially fails its readiness probe while the model weights load into the GPU. Traefik does not send any traffic to it. Meanwhile, Version 1 continues serving users. Only when Version 2 returns a 200 OK on its health endpoint does Traefik route traffic to it, and subsequently, Kubernetes begins draining and terminating the Version 1 pods.}
    \qa{What happens if the new model has a bug and never passes the readiness probe?}{The deployment halts. Version 2 remains in a pending state, and Traefik continues routing 100\% of traffic to Version 1 indefinitely. Users experience no disruption, and the engineering team is alerted to the failed deployment.}
\end{frame}
